% CORA poster abstract, ARC Prize Research Summit, 23 October 2026.
% Every number is indexed in records/CORA_SUMMIT_POSTER_EVIDENCE_20261012.md.
\documentclass[11pt]{article}
\usepackage[a4paper, margin=2.4cm]{geometry}
\usepackage[scaled]{helvet}
\renewcommand{\familydefault}{\sfdefault}
\usepackage[T1]{fontenc}
\usepackage{microtype}
\usepackage{xcolor}
\definecolor{ink2}{HTML}{52514E}
\pagestyle{empty}
\pdftrailerid{}
\setlength{\parskip}{0.55em}
\setlength{\parindent}{0pt}

\begin{document}
{\Large\bfseries Learning What to Change\par}
\vspace{0.2em}
{\large\bfseries Learning from Failure: Executable Reasoning and Hypothesis Revision\par}
\vspace{0.6em}
{\large Chimdi Walter Ndubuisi, Christian Fluharty, Pierce Taylor, Toni Kazic\par}
{\color{ink2} University of Missouri \quad$\cdot$\quad ARC Prize Research Summit, 23 October 2026 \quad$\cdot$\quad poster abstract\par}
\vspace{1.0em}

How can a machine learn the representation of new data, learn from failure, and know what to reconsider when its
explanation is not enough? We use ARC-AGI-2 as a controlled testbed: a few demonstrations define each task, and every
explanation can be executed and checked exactly.

CORA, our object-based program-synthesis system, writes explanations as executable rules over objects found in one
of seven segmentation views, runs them on every demonstration, compares rival rules, and certifies an answer only
when the rule, rebuilt with each demonstration hidden in turn, still predicts it. A separate evaluator scores every
study, under a protocol frozen before the run.

On 60 exposed development tasks, this check cut 8 right and 26 wrong first-ranked answers to 6, all right. Using its
own failure records, CORA built a guarded line-drawing rule for one real training task, rebuilt in 3 of 3 folds and
correct; a constructor that ignored the failure record found the same rule, so the value of the record is not
established. On 58 further training tasks the construction arms certified no new solution, and none wrong. A
version-space reasoner abstains when a rule entry it needs is supported by a single demonstration.

A trained guide on frozen Qwen3-4B features beat unguided search (198 vs 104 of 395 synthetic tasks) but not
hand-designed search (358). Rebuilding CORA's native program classes in a typed language raised held-out synthetic
pass@2 from 43 to 162 of 200, yet on 40 public training tasks it solved 2 (1 before), below the gain of 3 we required
in advance. Joint rule construction with evidence-aware selection then raised held-out synthetic pass@2 from 166 to
271 of 296 (synthetic tasks inside our own grammar: engineering progress, not ARC evidence). A descriptive analysis of
the real failures points at representation: which objects change is usually describable, but the changes often touch
parts of objects or add lines. We have built and tested rules over parts of objects, regions and lines; their frozen
study has not been run.

The limits are plain. CORA's Kaggle ARC-AGI-2 public score is 0.00\%; a reproduction of the published NVARC method,
which is not ours, scores 27.64\%. That failure records improve revision, and that a revision procedure transfers to
biological questions with independent validation, are both untested.

\vfill
{\small\color{ink2}
\textbf{Keywords:} ARC-AGI-2; program synthesis; learning from failure; hypothesis revision; representation;
verification; leave-one-out.\par
\textbf{Evidence tiers:} frozen synthetic studies (run once, designed grammar); development studies on exposed public
training tasks (not unseen scores); descriptive analyses after the fact; the official Kaggle public scores. Every
number, with its record and commit, is in \texttt{records/CORA\_SUMMIT\_CLAIM\_LEDGER\_20261017.md} (rows 1 to 31).\par}
\end{document}
